\chapter{Parsing expression grammars}

\section{Introdução}

No capítulo anterior, estudamos os combinadores de parsing como uma
forma de implementar analisadores sintáticos descendentes recursivos.
Neste capítulo, vamos estudar um formalismo que permite a definição
de analisadores sintáticos diretamente, sem a necessidade de usar uma
gramática livre de contexto. Esse formalismo são as \textbf{Parsing
Expression Grammars} (PEGs), introduzidas por Bryan Ford em 2004.

As PEGs compartilham com as gramáticas livres de contexto (GLCs) a
ideia de descrever linguagens por meio de regras de reescrita, mas
diferem em um ponto fundamental: enquanto as GLCs usam
\textbf{escolha não-determinística} entre as alternativas de uma
produção, as PEGs usam \textbf{escolha ordenada}. Em vez de tentar
todas as alternativas em paralelo (como faz um autômato
não-determinístico), uma PEG experimenta as alternativas em ordem e
se compromete com a primeira que tiver sucesso. Essa semântica
determinística elimina a ambiguidade presente em gramáticas e
expressões regulares.

O código discutido neste capítulo está disponível em
\begin{flushleft}
  \verb|PEG/Semantics/Simple.hs|.
\end{flushleft}

\section{Sintaxe e semântica de PEGs}

Uma PEG é composta por um conjunto de \textbf{regras de parsing},
cada uma associando um nome (não-terminal) a uma \textbf{expressão de
parsing}. As expressões de parsing são construídas a partir dos
seguintes operadores:

\begin{table}[ht]
  \centering
  \begin{tabular}{lll}
    \toprule
    Expressão & Nome & Descrição \\
    \midrule
    $\varepsilon$   & vazio              & Sempre tem sucesso sem
    consumir entrada \\
    $a$             & terminal           & Consome e casa com o símbolo $a$ \\
    $A$             & não-terminal       & Invoca a regra de nome $A$ \\
    $e_1\; e_2$     & sequência          & Tenta $e_1$ e depois $e_2$ \\
    $e_1 / e_2$     & escolha ordenada   & Tenta $e_1$; se falhar,
    tenta $e_2$ \\
    $e^*$           & repetição          & Zero ou mais repetições de
    $e$ (gulosa) \\
    $e^+$           & repetição positiva & Uma ou mais repetições de $e$ \\
    $e\,?$          & opcional           & Zero ou uma ocorrência de $e$ \\
    $!e$            & predicado negativo & Sucesso se $e$ falhar; não
    consome entrada \\
    $\& e$          & predicado positivo & Sucesso se $e$ tiver
    sucesso; não consome entrada \\
    \bottomrule
  \end{tabular}
  \caption{Operadores de PEG}
\end{table}

Uma PEG é então um conjunto de regras da forma \(A \leftarrow e\), em
que \(A\) é um não-terminal e \(e\) é uma expressão de parsing. Uma
das regras é designada como \textbf{regra inicial}.

Como exemplo, a seguinte PEG descreve a linguagem de expressões
aritméticas com precedência e associatividade já embutidas na gramática:

\[
  \begin{array}{lcl}
    E & \leftarrow & T\,(\texttt{+}\, T)^\star\\
    T & \leftarrow & F\,(*\, F)^\star\\
    F & \leftarrow & \texttt{(}\, E\, \texttt{)} \,/\, [0\text{-}9]^+
  \end{array}
\]

Comparada à GLC equivalente, a PEG não precisa de produções
recursivas à esquerda para expressar associatividade à esquerda: o
operador \((\cdot)^*\) captura diretamente a ideia de ``zero ou mais
ocorrências à direita'', que é processada iterativamente.

A semântica de uma PEG é definida em termos de uma função de parsing
que, dado uma expressão \(e\) e uma entrada \(s\), retorna ou um
\textbf{sucesso} que consiste de um par formado com a parte da
entrada consumida e o restante ainda não processado; ou uma \textbf{falha}.

A diferença crucial em relação às GLCs está na semântica da
\textbf{escolha ordenada} \(e_1 / e_2\):
Tenta-se \(e_1\) sobre a entrada \(s\). Se \(e_1\) \textbf{tem
sucesso}, o resultado é imediatamente retornado e \(e_2\) nunca é
tentada. Se \(e_1\) \textbf{falha}, então \(e_2\) é executada sobre
a entrada \(s\).

\subsection{Predicados}

Os predicados \(!e\) e \(\&\,e\) permitem analisar parte da entrada
sem necessariamente consumir parte dela. Temos os seguintes
operadores de predicado:

\begin{itemize}
  \item
    \(!e\) (\textbf{predicado negativo}): tem sucesso se e somente se
    \(e\) falha. Nenhuma entrada é consumida em nenhum caso.
  \item
    \(\&\,e\) (\textbf{predicado positivo}): tem sucesso se e somente
    se \(e\) tem sucesso. Nenhuma entrada é consumida. É equivalente
    a \(!(!e)\).
\end{itemize}

Os predicados são úteis para especificar restrições contextuais sem
avançar na entrada. Por exemplo, para reconhecer uma palavra
reservada como if sem aceitar um
identificador como iff, podemos escrever:

\[
  \texttt{if}\, !\,[a\text{-}zA\text{-}Z0\text{-}9]
\]

Isso garante que \haskell|if| só é reconhecido quando
não seguido de um caractere alfanumérico.

\subsection{Repetição}

O operador \(e^*\) é definido recursivamente como:

\[
  e^* \leftarrow e\, e^* \,/\\, \varepsilon
\]

Como a escolha é ordenada, a repetição é \textbf{gulosa}: o parser
sempre tenta consumir mais uma ocorrência de \(e\) antes de aceitar o
caso vazio. Isso garante que \(e^*\) sempre consumirá o maior prefixo
possível que case com \(e\).

Ao contrário das GLCs, não há ambiguidade em \(e^*\): o resultado é
determinístico e único para qualquer entrada.

\subsection{Terminação em PEGs}

PEGs podem ser vistas como uma representação de analisadores
descendentes recursivos e, portanto, possuem limitações similares. A
terminação do processo de parsing de um PEG é garantido de terminar
se a gramática atende as seguintes restrições:

\begin{itemize}
  \item
    Não possui regras recursivas à esquerda, diretas ou indiretas;
  \item
    Toda ocorrência de repetição, \(e^*\) é tal que \(e\) não deve
    ter sucesso sem consumir símbolos da entrada.
\end{itemize}

A primeira restrição é óbvia: pois faz com que analisadores
descendentes entrem em loop. A segunda é necessária devido à
semântica gulosa do operador de repetição: ele sempre tentará casar a
expressão \(e\) antes de finalizar a repetição. Se \(e\) tiver
sucesso sem consumir entrada, então \(e^*\) pode, a cada iteração,
consumir um prefixo vazio da entrada, o que leva a não terminação.

Dizemos que uma PEG é \textbf{bem formada} se ela atende as duas
restrições mencionadas. Para PEGs bem formadas há uma garantia formal
de que o processo de parsing sempre termina.

\section{Implementação em Haskell}

A implementação de PEGs em Haskell presente em
\verb|PEG/Semantics/Simple.hs| segue diretamente a
semântica apresentada. O ponto central é representar explicitamente
se o parser consumiu ou não parte da entrada, pois isso determina
quando a escolha ordenada pode retroceder.

\subsection{O tipo Result}

O resultado de um parser é representado pelo tipo:

\begin{minted}{haskell}
data Result d a
  = Pure a           -- sucesso sem consumir entrada
  | Commit d a       -- sucesso consumindo entrada; d: o resto da entrada
  | Fail String Bool -- falha; Bool indica se houve consumo
\end{minted}

Os três construtores capturam os três casos possíveis:

\begin{itemize}
  \item\haskell{Pure a}: o parser teve sucesso sem
    consumir nenhum símbolo. O resultado pode ser descartado e o
    parser alternativo pode ser tentado sobre a mesma entrada.
  \item\haskell{Commit d' a}: o parser teve sucesso e
    consumiu parte da entrada; \haskell{d'} é o
    restante. Falhas posteriores não causarão backtracking sobre este ponto.
  \item\haskell{Fail s c}: o parser falhou; a
    \haskell{String} \haskell{s}
    descreve o erro e o \haskell{Bool}
    \haskell{c} indica se houve consumo de entrada
    antes da falha.
\end{itemize}

\subsection{O tipo PExp}

Um parser PEG é um valor do tipo:

\begin{minted}{haskell}
newtype PExp d a
  = PExp { runPExp :: d -> Result d a }
\end{minted}

\haskell{PExp d a} representa um parser que lê de uma
entrada do tipo \haskell{d} (tipicamente \haskell{String})
e produz um resultado do tipo \haskell{a}. O tipo
\haskell{PExp} é uma instância de \haskell{Functor} que
permite modificar o valor de resultado do parser.

\subsection{Instâncias e combinadores}

A instância de \haskell{Applicative} define a
composição sequencial de dois parsers propagando o estado de consumo:

\begin{minted}{haskell}
instance Applicative (PExp d) where
  pure a = PExp $ \ _ -> Pure a
  PExp mf <*> PExp ma
    = PExp $ \ d ->
      case mf d of
        Pure f      -> fmap f (ma d)
        Fail s c    -> Fail s c
        Commit d' f ->
          case ma d' of
            Pure a       -> Commit d' (f a)
            Fail s _     -> Fail s True
            Commit d'' a -> Commit d'' (f a)
\end{minted}

O caso \haskell{Commit d' f} é o mais importante: se
o primeiro parser consumiu entrada, qualquer falha do segundo parser
é marcada como \haskell{Fail s True}, que indica que o consumo
ocorreu, sem possibilidade de retrocesso.

A instância de \haskell{Alternative} implementa a
semântica da \textbf{escolha}:

\begin{minted}{haskell}
instance Alternative (PExp d) where
  PExp ma <|> PExp mb
    = PExp $ \ d ->
      case ma d of
        Fail _ False -> mb d  -- falhou SEM consumir: tenta alternativa
        x            -> x     -- sucesso OU falhou COM consumo: retorna
  empty = PExp $ \ _ -> Fail "empty" False
\end{minted}

O padrão \haskell{Fail \_ False} é a chave: o segundo
parser \haskell{mb} só é tentado se o primeiro falhou
\textbf{sem consumir entrada} (\haskell{False}). Se o
primeiro consumiu entrada e falhou (\haskell{Fail \_ True}),
o erro é propagado imediatamente, sem retrocesso.

\subsection{O combinador try}

Para permitir retrocesso mesmo quando já houve consumo, a biblioteca
oferece o combinador \haskell{try}:

\begin{minted}{haskell}
try :: PExp d a -> PExp d a
try (PExp m)
  = PExp $ \d ->
      case m d of
        Fail s _ -> Fail s False  -- sempre marca como "sem consumo"
        x        -> x
\end{minted}

\haskell{try p} executa \haskell{p}
normalmente. Se \haskell{p} falha, com ou sem
consumo, \haskell{try} marca o resultado como
\haskell{Fail s False}, sinalizando que nenhum
consumo ocorreu para fins de retrocesso. Isso efetivamente
``cancela'' qualquer consumo parcial de \haskell{p},
permitindo que outras alternativas, caso existam, sejam tentadas.

\subsection{A escolha ordenada}

Combinando \haskell{try} e \haskell{<|>}, a biblioteca define
o operador de escolha ordenada:

\begin{minted}{haskell}
infixl 3 </>

(</>) :: PExp d a -> PExp d a -> PExp d a
p </> q = try p <|> q
\end{minted}

\haskell{p </> q} tenta \haskell{p}
com backtracking garantido: se \haskell{p} falha por
qualquer razão, \haskell{q} é tentado sobre a entrada original.

\subsection{O predicado not}

O combinador \haskell{not} implementa o predicado
negativo \(!e\):

\begin{minted}{haskell}
not :: PExp d a -> PExp d ()
not (PExp m)
  = PExp $ \d ->
      case m d of
        Fail{} -> Pure ()    -- e falhou: sucesso sem consumir
        _      -> Fail "unexpected" False  -- e teve sucesso: falha
\end{minted}

\haskell{not p} tem sucesso se e somente se
\haskell{p} falha, e nunca consome entrada. Sobre
ele, \haskell{eof} é definido como:

\begin{minted}{haskell}
eof :: Stream d => PExp d ()
eof = not anyChar
\end{minted}

\haskell{eof} tem sucesso quando não há mais entrada disponível.

\subsection{Combinadores léxicos}

A biblioteca oferece combinadores para reconhecimento de tokens comuns:

\begin{minted}{haskell}
satisfy :: Stream d => (Char -> Bool) -> PExp d Char
satisfy p = try $ do
  x <- anyChar
  x <$ guard (p x)

char :: Stream d => Char -> PExp d Char
char c = satisfy (c ==)

lexeme :: Stream d => PExp d a -> PExp d a
lexeme m = m <* whiteSpace

symbol :: Stream d => Char -> PExp d Char
symbol c = lexeme (char c)

digit :: Stream d => PExp d Int
digit = f <$> satisfy isDigit
  where f c = ord c - ord '0'
\end{minted}

O combinador \haskell{satisfy} usa
\haskell{try} para garantir que, se o predicado falha
(o caractere foi consumido por \haskell{anyChar} mas
não satisfaz \haskell{p}), o parser retrocede e marca
a falha como sem consumo. Isso é necessário para que
\haskell{satisfy} se comporte corretamente como
primitivo: uma falha em \haskell{satisfy} não deve
comprometer a entrada.

O combinador \haskell{lexeme} descarta espaços em
branco após o item reconhecido, e \haskell{symbol c}
combina o reconhecimento de um caractere específico com descarte de
espaço, o mesmo padrão que vimos no Megaparsec.

\section{Exemplo: parser de expressões}

Para ilustrar o uso da biblioteca, vamos construir um parser de
expressões aritméticas. A gramática PEG é:

\[
  \begin{array}{lcl}
    E & \leftarrow & T\, (\texttt{+}\, T)^\star\\
    T & \leftarrow & F\,(*\, F)^\star\\
    F & \leftarrow & \texttt{(}\, E\, \texttt{)} \,/\, [0\text{-}9]^+
  \end{array}
\]

Em Haskell, usando os combinadores da biblioteca:

\begin{minted}{haskell}
data Exp = Lit Int | Exp :+: Exp | Exp :*: Exp

expP :: Stream d => PExp d Exp
expP = foldl (:+:) <$> termP <*> many (symbol '+' *> termP)

termP :: Stream d => PExp d Exp
termP = foldl (:*:) <$> factP <*> many (symbol '*' *> factP)

factP :: Stream d => PExp d Exp
factP = between (symbol '(') (symbol ')') expP
    </> Lit . foldl (\a b -> a * 10 + b) 0 <$> many1 digit
\end{minted}

A estrutura do parser segue diretamente a gramática PEG. Em
\haskell{factP}, o operador
\haskell{</>} garante que, se a tentativa de
reconhecer uma expressão com parêntesis falhar, o parser retroceda e
tente reconhecer um literal inteiro.

Note que \haskell{expP} e
\haskell{termP} usam \haskell{many}
com \haskell{foldl} para construir a AST com
associatividade à esquerda, sem precisar de recursão à esquerda nas
regras. Isso é uma vantagem das PEGs sobre as GLCs: associatividade à
esquerda é expressa naturalmente por repetição, sem introduzir
recursão problemática para parsers descendentes.

\section{PEGs e GLCs: principais diferenças}

As PEGs e as GLCs são formalismos distintos para descrever
linguagens. As principais diferenças práticas são:

\textbf{Ambiguidade.} Toda PEG é, por definição, não-ambígua: a
escolha ordenada garante que exista no máximo um resultado para
qualquer entrada. As GLCs podem ser ambíguas, e resolver ambiguidades
muitas vezes exige reformular a gramática ou usar anotações de
precedência externas (como os mecanismos
  \verb|%left|/\verb|%right| do yacc/bison).

  \textbf{Recursão à esquerda.} GLCs admitem recursão à esquerda
  naturalmente; por exemplo, \(E \to E + T\) é uma produção válida e
  útil. Parsers PEG descendentes não suportam recursão à esquerda
  direta, ela causaria um laço infinito. A solução é reescrever
  usando repetição \((\cdot)^*\) ou usar técnicas especializadas de
  eliminação de recursão à esquerda.

  \textbf{Poder expressivo.} PEGs reconhecem linguagens que as GLCs não
  conseguem (por exemplo, \(\{a^n b^n c^n \mid n \geq 0\}\), usando
  predicados) e há conjectura de que algumas linguagens livres de
  contexto não podem ser descritas por PEGs, embora isso não esteja
  formalmente provado.

  \textbf{Mensagens de erro.} A semântica de PEGs facilita a geração de
  mensagens de erro precisas: quando o parser se compromete com uma
  alternativa e depois falha, sabe exatamente onde e o que era esperado.

  \section{Conclusão}

  Este capítulo apresentou as Parsing Expression Grammars como um
  formalismo para especificar analisadores sintáticos determinísticos.
  As PEGs se distinguem das GLCs pela \textbf{escolha ordenada} \(/\),
  que elimina a ambiguidade e corresponde diretamente ao
  comportamento de parsers.

  Apresentamos a implementação em Haskell da biblioteca
  \verb|PEG.Semantics.Simple|, cujo tipo
  \haskell{Result d a} representa explicitamente os
  três estados possíveis de um parser: sucesso sem consumo
  (\haskell{Pure}), sucesso com consumo
  (\haskell{Commit}) e falha (\haskell{Fail}).

  Os combinadores \haskell{<|>} (escolha sem
  backtracking) e \haskell{</>} (escolha com
  backtracking via \haskell{try}) dão ao programador
  controle explícito sobre quando o parser pode recuar, tornando as
  PEGs uma ferramenta poderosa tanto para especificação formal quanto
  para implementação prática de analisadores sintáticos.
